\section{Atmosfera bosimi}

\begin{enumerate}
\item Torrichelli tajribasi uzunliklari 1 m dan bo‘lgan, ko‘ndalang kesimlari 
farq qiluvchi 3 ta nayda o‘tkazilsa, naylardagi simob ustunlarining 
balandliklari farq qiladimi?

\begin{enumerate}[label=\Alph*)]
\item $h_1>h_2>h_3$
\item $h_1<h_2>h_3$
\item $h_1>h_2<h_3$
\item $h_1=h_2=h_3$
\end{enumerate}



\item Torichelli tajribasida ishlatiladigan simobli barometr 
(bir uchi berk shisha nay va simobli idish) biror balandlikdan 
vertikal tashlab yuborilsa, nay ichidagi simob sathining balandligi 
qanday o‘zgaradi?

\begin{enumerate}[label=\Alph*)]
\item o‘zgarmaydi
\item simob nayni butunlay to‘ldiradi
\item nay ichidagi simob butunlay oqib chiqadi
\item avval oshadi, so‘ngra o‘zgarmaydi
\end{enumerate}


\item Havosi so‘rib olingan nayni normal sharoitda simobga tushirilsa, 
simob atmosfera bosimi ta’sirida shu nay orqali 76 cm ko‘tariladi. 
Shu nay suvga tushirilsa, suv necha metr balandlikka ko‘tarilgan bo‘lar edi?

\begin{enumerate}[label=\Alph*)]
\item 10,33
\item 8,5
\item 9,8
\item 11
\end{enumerate}


\item Diametrlari $d_1<d_2<d_3<d_4$ munosabatda bo‘lgan bir uchi berk 
4 ta shisha naycha 4 xil suyuqlik bilan to‘ldirilib, ochiq uchi bilan 
o‘z suyuqliklariga to‘nkarilganda, ulardagi suyuqliklarning balandliklari 
$h_1=h_4<h_2=h_3$ munosabatda bo‘ladi. 
Suyuqlik zichliklari uchun qaysi munosabat o‘rinli?

\begin{enumerate}[label=\Alph*)]
\item $\rho_4>\rho_3>\rho_2>\rho_1$
\item $\rho_4>\rho_1>\rho_2>\rho_3$
\item $\rho_1=\rho_4>\rho_2=\rho_3$
\item $\rho_2>\rho_3,\ \rho_1>\rho_4$
\end{enumerate}


\item Toshkent shahri dengiz sathidan 407 m balandlikda joylashgan. 
Yunusoboddagi televizion minoraning balandligi 385 m. 
Agar dengiz sathidagi atmosfera bosimi 760 mm Hg ga teng bo‘lsa, 
minoraning uchidagi barometr qanday bosimni (mm Hg) ko‘rsatadi?

\begin{enumerate}[label=\Alph*)]
\item 704
\item 698
\item 694
\item 684

\end{enumerate}


\item Shaxta tubida atmosfera bosimi 820 mm simob ustuniga teng bo‘lsa, 
shaxta chuqurligi necha metr?

\begin{enumerate}[label=\Alph*)]
\item 820
\item 760
\item 600
\item 720
\end{enumerate}


\item Toshkent shahri dengiz sathidan 408 m balandlikda. 
Televizion minoraning balandligi 192 m. 
Minoraning tepasida barometr necha mm Hg bosimni ko‘rsatadi? 
100 m ga ko‘tarilganda bosim 9 mm Hg ga kamayadi deb hisoblang. 
Dengiz sathidagi bosim 760 mm Hg.

\begin{enumerate}[label=\Alph*)]
\item 760
\item 706
\item 728
\item 710

\end{enumerate}

\end{enumerate}



\begin{table}[h!]
\centering
\begin{tabular}{|c|c|c|c|c|c|c|}
\hline
№ & J. & № & J. & № & J.  \\ \hline
1 & D & 6 & D & 11 &  \\ \hline
2 & A & 7 & B & 12 &  \\ \hline
3 & A & 8 &  & 13 &  \\ \hline
4 & C & 9 &  & 14 &   \\ \hline
5 & B & 10 &  & 15 &   \\ \hline
\end{tabular}
\end{table}